\documentclass[12pt]{article}
\usepackage[T1]{fontenc}
\usepackage[utf8]{inputenc}
\usepackage{lmodern}
\usepackage{textcomp}
\usepackage{enumitem}
\usepackage{geometry}
\geometry{margin=1in}
\begin{document}
\begin{center}
Answer Key - Astronomy - Graduate Level Assessment 24 \
\small (Answers are ordered exactly as in the question paper.)
\end{center}

\section*{Multiple Choice}
\begin{enumerate}[label=\arabic*.]
\item \textbf{Answer:} C
\\ \textit{Options:} A) It has a fusion crust; B) It contains solidified spherical droplets; C) It is highly processed; D) It has different elemental composition than earth
\\ \textit{Note:} A highly processed rock typically indicates alteration by geological processes on Earth, which is not a characteristic of meteorites that generally retain their primordial state from space.
\item \textbf{Answer:} C
\\ \textit{Options:} A) A kilogram of feathers; B) Five pounds of bricks as measured on Earth; C) Five kilograms of feathers; D) A kilogram of bricks
\\ \textit{Note:} Five kilograms of feathers would weigh the most on the moon because weight is a measure of mass affected by the gravitational pull, and five kilograms is the largest mass among the options presented, regardless of the material's composition.
\item \textbf{Answer:} A
\\ \textit{Options:} A) Its rotation is too slow.; B) It has too thick an atmosphere.; C) It is too large.; D) It does not have a metallic core.
\\ \textit{Note:} Venus does not have a strong magnetic field primarily because its slow rotation, which takes about 243 Earth days, is insufficient to generate the dynamo effect necessary for creating a substantial magnetic field like that of Earth.
\item \textbf{Answer:} C
\\ \textit{Options:} A) viruses; B) bacteria such as E. coli; C) organisms living deep in the oceans around seafloor volcanic vents and in hot springs; D) plankton that use sunlight as an energy source through photosynthesis
\\ \textit{Note:} Organisms living deep in the oceans around seafloor volcanic vents and in hot springs are considered to most resemble the common ancestor of all life because they are extremophiles that thrive in harsh environments similar to those of early Earth, exhibiting ancient metabolic pathways and genetic traits that provide insight into the origins of life.
\item \textbf{Answer:} B
\\ \textit{Options:} A) Being hit in the head by a bullet.; B) Being hit by a small meteorite.; C) Starvation during global winter caused by a major impact.; D) Driving while intoxicated without wearing seatbelts.
\\ \textit{Note:} Being hit by a small meteorite is the least likely cause of death because such occurrences are extremely rare, with very few documented cases in human history, making the probability of fatality from this event negligible compared to other causes of death.
\end{enumerate}

\section*{True / False}
\begin{enumerate}[label=\arabic*.]
\setcounter{enumi}{5}
\item \textbf{Answer:} False
\\ \textit{Options:} A) True; B) False
\\ \textit{Note:} A lunar day, which is the time it takes for the Moon to complete one rotation on its axis relative to the Sun, lasts approximately 29.5 Earth days, significantly longer than the 24-hour day on Earth.
\item \textbf{Answer:} False
\\ \textit{Options:} A) True; B) False
\\ \textit{Note:} Callisto is the third largest of Jupiter's four largest moons, orbiting at a greater distance from Jupiter than Io, which is the closest of the four.
\item \textbf{Answer:} False
\\ \textit{Options:} A) True; B) False
\\ \textit{Note:} Kirkwood gaps occur at specific orbital resonances with Jupiter, particularly where asteroids have orbital periods that are in a simple ratio with Jupiter's, not Mars.
\item \textbf{Answer:} True
\\ \textit{Options:} A) True; B) False
\\ \textit{Note:} The statement is true because when light from a hot star passes through cooler hydrogen gas, specific wavelengths of light are absorbed by hydrogen atoms, resulting in dark absorption lines in the star's continuous spectrum.
\item \textbf{Answer:} True
\\ \textit{Options:} A) True; B) False
\\ \textit{Note:} The statement is true because the International Space Station orbits Earth at an average altitude of about 400 kilometers, requiring it to travel at a speed of approximately 28,000 km/h to maintain a stable low Earth orbit and counteract gravitational pull.
\end{enumerate}

\section*{Long Form Response}
\begin{enumerate}[label=\arabic*.]
\setcounter{enumi}{10}
\item \textbf{Answer:} Rapid rotation plays a crucial role in the formation of the equatorial bulge observed in jovian planets, such as Jupiter and Saturn. As these planets rotate at high speeds, the centrifugal force generated acts to fling mass outward, particularly at the equator, resulting in a noticeable bulge. This phenomenon is primarily a consequence of the balance between gravitational forces, which pull the mass inward, and the outward centrifugal force generated by rotation. The implications for their structure are significant, as this distribution of mass not only affects the planets' shapes but also influences their gravitational fields and atmospheric dynamics. Moreover, the equatorial bulge can impact the planets' rotational characteristics and their satellite systems, underscoring the intricate interplay between rotation and planetary structure in the context of astrophysical phenomena.
\\ \textit{Note:} correct because it accurately describes how the interplay between gravitational forces and centrifugal force due to rapid rotation leads to the formation of an equatorial bulge in jovian planets, highlighting the physical mechanisms at work and their structural implications.
\item \textbf{Answer:} The angular resolution of the human eye, typically around 1 arcminute, significantly influences the ability to distinguish closely spaced stars. When two stars are separated by an angular distance smaller than this threshold, they appear as a single point of light rather than two distinct sources. This limitation not only affects amateur observations but also has implications for astronomical studies, particularly in identifying binary star systems or resolving crowded fields in star clusters. As a result, the inability to differentiate between closely spaced stars can lead to misinterpretations of stellar properties and dynamics, emphasizing the need for telescopes with higher resolution capabilities to enhance our understanding of the universe.
\\ \textit{Note:} correct because it accurately describes the angular resolution of the human eye as a limiting factor in distinguishing closely spaced stars, which impacts both amateur observations and professional astronomical studies, particularly in identifying binary stars and resolving star clusters.
\end{enumerate}

\vfill
\noindent\textit{End of Answer Key}
\end{document}